
%
%  $Description: Author guidelines and sample document in LaTeX 2.09$ 
%
%  $Author: ienne $
%  $Date: 1995/09/15 15:20:59 $
%  $Revision: 1.4 $
%

\documentclass[times, 10pt,twocolumn]{article} 
\usepackage{latex8}
\usepackage{times}
\usepackage{enumitem}
\usepackage{graphicx}
\usepackage{booktabs}
\usepackage{url}
\setitemize{noitemsep,topsep=0pt,parsep=0pt,partopsep=0pt}
\input{header}
%\documentstyle[times,art10,twocolumn,latex8]{article}

%------------------------------------------------------------------------- 
% take the % away on next line to produce the final camera-ready version 
\pagestyle{empty}
%------------------------------------------------------------------------- 
\begin{document}
\title{Causal Metadata Growth Under Membership Churn: A Simulation Study of Vector Clocks and Dotted Version Vectors}

\author{
Eric Nelson$^{*}$ \\
Department of Computer Science\\
The University of Texas at Austin \\
Austin, Texas, US \\
eric.nelson.swe@gmail.com
\and
Pallavi Desai$^{*}$ \\
Department of Computer Science\\
The University of Texas at Austin \\
Austin, Texas, US \\
pallavi.desai@utexas.edu
}
\maketitle
\thispagestyle{empty}
\begingroup
\renewcommand\thefootnote{*}
\footnotetext{Equal contribution}
\endgroup

\begin{abstract}
  Here is the abstract. 
  Code is available at: \url{https://github.com/pallavidesai-ut/distributed-computing-project/}
\end{abstract}
\vspace{-2em}
\Section{Core Idea}
\vspace{-1em}

In distributed systems that prioritize availability and partition tolerance, nodes must make progress independently and reconcile state asynchronously. Causal consistency has emerged as a principled middle ground between strong consistency and eventual consistency: it preserves the intuitive "happens-before" relationship between operations while allowing concurrent writes to proceed without coordination. The standard mechanism for tracking causality is the vector clock, which assigns each node a logical counter and requires every message to carry a snapshot of the entire counter vector. This approach is correct and well-understood, but it carries a fundamental scalability problem: the metadata attached to every message grows linearly with the number of distinct nodes that have ever participated in the system.

Dotted Version Vectors (DVVs) were proposed as a more precise alternative. Rather than treating the clock as a single monolithic counter vector, DVVs separate the dot, the specific write event being described , from the context, the causal history it depends on. This separation allows the system to distinguish concurrent writes from duplicate deliveries without requiring a full per-node history, and it creates a natural opportunity for pruning: once a dot is no longer needed to resolve concurrency, it can be retired. An extension to this model introduces lease-based expiration, where entries for departed nodes are removed after a configurable time window, accepting a bounded risk of causal violation in exchange for meaningful metadata reduction.

\vspace{-1em}

\Section{Motivation}
\vspace{-1em}

The metadata growth problem is especially acute in dynamic deployments where nodes join and leave frequently. In a stable cluster, vector clock entries are bounded by the fixed membership size. Under churn, however, departed nodes leave behind dead entries that continue to propagate with every message. No standard pruning mechanism exists, because removing an entry risks silently breaking causal delivery guarantees for any node that has not yet seen the messages that entry tracks. The result is unbounded metadata bloat that inflates message size, increases buffering overhead, and degrades system performance even after the nodes responsible for that metadata have long since left.

Despite the theoretical appeal of DVVs, it remains unclear how much practical benefit they provide under realistic churn conditions, what the actual tradeoff curve looks like between metadata reduction and causal safety when leases are applied, and whether the difference is large enough to matter for latency and buffering in a real workload. These questions motivate the present work.

We implement a discrete-event simulation of a distributed key-value store in which multiple nodes broadcast writes to their peers, maintain causal ordering metadata, and buffer messages whose dependencies have not yet been satisfied. The simulation supports four membership churn profiles (stable, low, sustained, and burst) and is designed to run the same workload across three clock implementations: a standard VectorClock, a DVV, and a lease-based DVV with configurable expiration windows. We measure metadata size per message, delivery latency, buffering pressure, and causal violation rate across all combinations, allowing us to directly compare the cost and correctness properties of each approach.

Our central research question is: does causal metadata size grow with membership churn rate, and can Dotted Version Vectors with an optional lease-based pruning extension reduce that growth compared to standard vector clocks without sacrificing causal delivery guarantees? A secondary question concerns the lease-based extension specifically: what is the tradeoff curve between metadata reduction and causal violation rate as the lease window shrinks?

\vspace{-1em}

\Section{Design}
\vspace{-1em}

Our simulation models a distributed key-value store as a collection of nodes that periodically broadcast writes to all active peers. Each write picks a random key and value, attaches causal metadata from the node's local clock, and sends a copy to every other active node. On receipt, a node checks whether the message's causal dependencies are satisfied before delivering it. If they are not, the message enters a buffer and is retried whenever a subsequent delivery updates the node's clock state. This models the core buffering pressure that arises when messages arrive out of causal order.

The simulation is driven by a discrete-event engine that schedules callbacks on a heap queue and processes them in time order. Membership churn is layered on top: a Cluster object manages the active node set and fires join and leave events according to one of four named profiles. The stable profile holds membership fixed throughout the run. The low profile introduces a slow trickle of joins and departures. The sustained profile applies continuous moderate churn. The burst profile periodically removes a large fraction of nodes and then reintroduces them, stress-testing how clocks behave when membership changes dramatically and repeatedly.

All three clock implementations (VectorClock, DVV, and LeaseBasedDVV) conform to a shared interface that exposes four operations: incrementing on a local write, building the metadata payload for an outgoing message, checking and updating state on receipt, and reporting the current metadata size. This design keeps the simulation harness clock-agnostic; swapping implementations requires no changes to node or cluster logic. A NullClock baseline, which attaches no metadata and enforces no ordering, provides a lower bound on all measurements. Metrics are collected continuously across each run and summarized as average metadata size per message, average delivery latency, total buffered messages, and causal violation count, which are the four quantities needed to evaluate both the efficiency and correctness of each clock under each churn profile.

\Section{Experimental Setup and Results}
\vspace{-1em}

Experimental results here. 

\vspace{-1em}
\Section{Conclusion}
\vspace{-1em}

Conclusion here. 

\vspace{-1em}
\bibliographystyle{latex8}
\bibliography{latex8}
\end{document}

